% ---------------------------------------------------------------------
\chapter{Gute Commits}
\label{kap:commits}

\begin{lernziele}
\begin{itemize}[nosep]
  \item wissen, wie gross ein Commit sein sollte und warum,
  \item Commit-Nachrichten nach gängiger Industrie-Konvention schreiben,
  \item die Konvention \emph{Conventional Commits} kennen.
\end{itemize}
\end{lernziele}

\section{Warum die Qualität von Commits zählt}

Ein Commit wird einmal geschrieben, aber oft gelesen: beim Review, bei der Fehlersuche
Monate später, von neuen Teammitgliedern. Die Geschichte eines Projekts ist seine
Dokumentation über das \emph{Warum}. Der Code zeigt, \emph{was} das Programm tut. Nur die
Commits erklären, \emph{weshalb} es so geworden ist.

\section{Ein Anliegen pro Commit}

\begin{definition}{Atomarer Commit}
Ein \begriff{atomarer Commit} enthält genau ein abgeschlossenes Anliegen, zum Beispiel eine
Funktion, eine Fehlerbehebung oder eine Umbenennung. Er lässt sich für sich allein
verstehen, prüfen und bei Bedarf rückgängig machen.
\end{definition}

\begin{center}
\renewcommand{\arraystretch}{1.3}
\begin{tabularx}{\linewidth}{@{}>{\raggedright\arraybackslash}X>{\raggedright\arraybackslash}X@{}}
  \toprule
  \textbf{Ungünstig} & \textbf{Besser} \\
  \midrule
  ein Commit „Abend 3“ mit Stack, Queue, Tippfehler in der README und neuer Ordnerstruktur
  & vier Commits: Stack, Queue, README korrigiert, Ordnerstruktur \\
  ein Commit pro gespeicherter Datei, auch halbfertig
  & ein Commit, sobald ein Teil funktioniert \\
  \bottomrule
\end{tabularx}
\end{center}

Stellt sich eine Woche später heraus, dass die Queue fehlerhaft ist, lässt sich im
zweiten Fall genau dieser eine Commit rückgängig machen, ohne Stack und README anzurühren.

\section{Die Commit-Nachricht}

Über Jahre hat sich ein Format etabliert, das von Werkzeugen wie GitHub und \cmd{git log}
optimal dargestellt wird:

\begin{lstlisting}[style=ausgabe]
ADS U02: Stack mit verketteter Liste implementiert

Der Array-Stack musste bei jeder Vergrösserung umkopieren. Die
verkettete Liste wächst pro Element in konstanter Zeit, was bei
unbekannter Datenmenge planbarer ist.
\end{lstlisting}

\begin{description}[style=nextline, leftmargin=1.2cm]
  \item[Betreffzeile] höchstens etwa 50 bis 70 Zeichen, beschreibt \emph{was} der Commit
    bewirkt. Sie erscheint in \cmd{git log --oneline} und als Titel auf GitHub.
  \item[Leerzeile] trennt Betreff und Text. Werkzeuge erkennen daran die Betreffzeile.
  \item[Text (optional)] erklärt \emph{warum}, wenn es nicht offensichtlich ist.
    Zeilen nach etwa 72 Zeichen umbrechen.
\end{description}

Mehrzeilige Nachrichten lassen sich auch ohne Editor schreiben: Jedes weitere \cmd{-m}
wird zu einem eigenen Absatz.

\begin{lstlisting}[style=befehle]
git commit -m "ADS U02: Stack mit verketteter Liste implementiert" -m "Der Array-Stack ..."
\end{lstlisting}

\subsection*{Gute und schlechte Beispiele}

\begin{center}
\renewcommand{\arraystretch}{1.3}
\begin{tabularx}{\linewidth}{@{}>{\ttfamily}l>{\raggedright\arraybackslash}X@{}}
  \toprule
  \normalfont\textbf{Nachricht} & \textbf{Beurteilung} \\
  \midrule
  update            & sagt nichts; was wurde aktualisiert, und warum? \\
  fix               & welcher Fehler? \\
  asdf              & leider häufig, völlig wertlos \\
  Änderungen        & beschreibt jeden denkbaren Commit \\
  \midrule
  LIAL SW03: LR-Zerlegung mit Pivotsuche     & klar, auffindbar, mit Bereich \\
  Spickzettel ADS: Heap-Operationen ergänzt  & klar, was und wo \\
  README: Einrichtung für Windows beschrieben & klar \\
  \bottomrule
\end{tabularx}
\end{center}

\section{Conventional Commits}

Viele Open-Source-Projekte und Firmen verwenden die Konvention \begriff{Conventional
Commits}. Jede Betreffzeile beginnt mit einem Typ:

\begin{center}
\renewcommand{\arraystretch}{1.25}
\begin{tabularx}{\linewidth}{@{}>{\ttfamily}l>{\raggedright\arraybackslash}X>{\ttfamily}l@{}}
  \toprule
  \normalfont\textbf{Typ} & \textbf{Bedeutung} & \normalfont\textbf{Beispiel} \\
  \midrule
  feat:     & neue Funktion          & feat: Export als CSV \\
  fix:      & Fehlerbehebung         & fix: Absturz bei leerer Liste \\
  docs:     & nur Dokumentation      & docs: Git-Handbuch Kapitel 9 \\
  refactor: & Umbau ohne neues Verhalten & refactor: Stack in eigene Klasse \\
  test:     & nur Tests              & test: Randfälle für Heap \\
  chore:    & Wartung, Konfiguration & chore: .gitignore ergänzt \\
  \bottomrule
\end{tabularx}
\end{center}

\begin{praxis}
Der Typ ist mehr als Ordnung: Werkzeuge lesen ihn automatisch aus. Aus allen
\cmd{feat:}- und \cmd{fix:}-Commits entstehen so die Versionshinweise einer neuen
Software-Version, ohne dass jemand sie von Hand schreibt.
\end{praxis}

\begin{erfahrung}
Ein guter Test für jede Commit-Nachricht: Ergibt der Satz „Wenn man diesen Commit
übernimmt, dann wird \ldots“ mit der Betreffzeile einen sinnvollen Satz? „\ldots{} dann wird
der Stack mit verketteter Liste implementiert“: ja. „\ldots{} dann wird update“: nein.
\end{erfahrung}

\begin{zusammenfassung}
\begin{itemize}[nosep]
  \item Ein Commit, ein Anliegen. Kleine Commits lassen sich prüfen und gezielt rückgängig machen.
  \item Betreffzeile kurz und konkret, Leerzeile, dann bei Bedarf das Warum.
  \item Conventional Commits (\cmd{feat:}, \cmd{fix:}, \ldots) ist ein verbreiteter Standard.
\end{itemize}
\end{zusammenfassung}
